\documentclass[12pt]{article}
\usepackage[T1]{fontenc}
\usepackage[utf8]{inputenc}
\usepackage{lmodern}
\usepackage{textcomp}
\usepackage{enumitem}
\usepackage{geometry}
\geometry{margin=1in}
\begin{document}
\begin{center}
Answer Key - Constitutional Law - Graduate Level Assessment 16 \
\small (Answers are ordered exactly as in the question paper.)
\end{center}

\section*{Multiple Choice}
\begin{enumerate}[label=\arabic*.]
\item \textbf{Answer:} A
\\ \textit{Options:} A) The Nazi law in question was validly enacted.; B) The Nazi rule of recognition was unclear.; C) Hart agreed with Fuller's argument.; D) The court misunderstood the legislation.; E) Fuller misconstrued the purpose of the law.
\\ \textit{Note:} Hart's response emphasizes that validity is determined by the procedures through which laws are enacted, and since Nazi laws followed the established legal processes of their time, they were valid despite their moral shortcomings.
\item \textbf{Answer:} E
\\ \textit{Options:} A) Such conduct is a public act but is shielded by immunity; B) Such conduct is a private act but is not shielded by immunity; C) Such conduct is a public act and does not attract immunity; D) Such conduct is a public act (jure imperii); E) Such conduct is a private act (jure gestionis)
\\ \textit{Note:} The unlawful homicide committed by the Minister of country X abroad is considered an act jure gestionis because it involves private conduct that does not fall within the sovereign functions of the state, distinguishing it from acts jure imperii, which pertain to governmental or sovereign acts.
\item \textbf{Answer:} C
\\ \textit{Options:} A) Purely discussed jurisprudence only; B) Defined law in according with morality and purity; C) Separated law from religion, ethics, sociology and history; D) Discussed law purely in terms of justice
\\ \textit{Note:} Kelsen's pure theory of law is deemed "pure" because it focuses exclusively on the normative structure of law, intentionally excluding influences from religion, ethics, sociology, and history to establish a clear and objective understanding of legal norms.
\item \textbf{Answer:} C
\\ \textit{Options:} A) Prior screening by neutral fact finder; B) Prior screening, notice and opportunity to respond; C) Prior notice and prior evidentiary hearing; D) Prior notice and immediate termination; E) Immediate termination without notice
\\ \textit{Note:} The correct answer is based on the principle of due process, which mandates that individuals must receive prior notice and have the opportunity for a hearing before the state can terminate parental custody rights or welfare benefits, ensuring fair treatment under the law.
\item \textbf{Answer:} A
\\ \textit{Options:} A) It distinguishes social rules from mere group habits.; B) It defines the judicial function.; C) It illustrates the authority of the legislature.; D) It stresses the relationship between law and justice.
\\ \textit{Note:} The 'internal point of view' in Hart's concept of law is crucial because it highlights the normative aspect of social rules, differentiating them from mere habits by emphasizing the participants' recognition and acceptance of these rules as standards for behavior.
\end{enumerate}

\section*{True / False}
\begin{enumerate}[label=\arabic*.]
\setcounter{enumi}{5}
\item \textbf{Answer:} True
\\ \textit{Options:} A) True; B) False
\\ \textit{Note:} Defining the limits of government power is typically regarded as a secondary function of the law within constitutional law frameworks, as the primary focus is on establishing rights and liberties for individuals.
\item \textbf{Answer:} True
\\ \textit{Options:} A) True; B) False
\\ \textit{Note:} The statement is true because Austin's theory of law emphasizes that laws are commands issued by a sovereign authority and are backed by the threat of sanctions or legal remedies for non-compliance, thus encapsulating the essential characteristics of command, sovereignty, and remedial provisions in legal systems.
\item \textbf{Answer:} False
\\ \textit{Options:} A) True; B) False
\\ \textit{Note:} Formalism emphasizes adherence to established legal rules and principles, minimizing the role of moral considerations in judicial decision-making.
\item \textbf{Answer:} True
\\ \textit{Options:} A) True; B) False
\\ \textit{Note:} Contracts for the sale of illegal substances are considered void and unenforceable because they violate public policy and the law, thereby lacking the legal capacity to form a binding agreement.
\item \textbf{Answer:} False
\\ \textit{Options:} A) True; B) False
\\ \textit{Note:} Aristotle argued that natural law is based on human nature and reason, rather than being solely derived from divine commands, emphasizing the role of human rationality in understanding moral principles.
\end{enumerate}

\section*{Long Form Response}
\begin{enumerate}[label=\arabic*.]
\setcounter{enumi}{10}
\item \textbf{Answer:} Catharine MacKinnon argues that the power dynamics between men and women fundamentally shape societal structures, positing that these dynamics are entrenched in law and culture, resulting in systemic inequality. She emphasizes that power is not merely a matter of individual strength but is deeply influenced by gendered relationships and societal norms that privilege male authority. In contrast, the perspective that power arises solely from individual strength overlooks the institutionalized nature of gender discrimination and fails to account for how societal structures perpetuate male dominance and female subordination. This individualistic view diminishes the significance of the collective power imbalances that MacKinnon highlights, which are crucial for understanding the systemic oppression faced by women in various spheres of life, including legal contexts. Ultimately, recognizing the intricate interplay between power and gender is essential for addressing inequality in constitutional law and society as a whole.
\\ \textit{Note:} correct because it highlights MacKinnon's argument that societal structures are shaped by entrenched power dynamics between genders, emphasizing that power is influenced by institutional and cultural factors rather than just individual strength.
\item \textbf{Answer:} Hobbes and Locke present fundamentally different views of life prior to the social contract, which significantly influences their theories of government. Hobbes characterizes the state of nature as a chaotic environment marked by fear and violence, where life is "solitary, poor, nasty, brutish, and short." This bleak perspective leads him to advocate for a powerful, centralized authority to maintain order and prevent anarchy. In contrast, Locke envisions the state of nature as a generally peaceful realm governed by natural rights and reason, where individuals possess inherent rights to life, liberty, and property. Consequently, Locke's theory promotes a government that is limited and accountable to the governed, emphasizing the protection of these natural rights. Thus, the contrasting depictions of pre-contractual life shape Hobbes' advocacy for absolute sovereignty and Locke's support for democratic governance based on consent.
\\ \textit{Note:} correct because it accurately contrasts Hobbes' pessimistic view of the state of nature, which necessitates a strong government for order, with Locke's more optimistic perspective that allows for a government based on consent and the protection of natural rights.
\end{enumerate}

\vfill
\noindent\textit{End of Answer Key}
\end{document}